\scanpage{9}

de façon générale, \deleted{$D_\xi(\alpha)$} $\sigma_i(\alpha)$ est une forme de degré $2i$ en $\alpha$, On a en particulier, par un calcul \unsure{facile}
\[
\mathrm{tr}(\alpha)=\frac{r}{I(\xi,\cdots,\xi)}I(\xi,\cdots,\xi,D(\alpha))
\]

\textcircled{d} Pour tout $A\in\mathrm{Ob}\,\unsure{\mathscr A|\mathcal A}$, on se donne deux parties
\[
N^{\unsure{>|+|\geqslant}}\ \deleted{\text{\unsure{et}}\ N^{\unsure{+|>}}}(A)\subset N(A)
\]
\deleted{$N^+(A)$ \unsure{tels} dépendent \unsure{de} \unsure{fonctoriellement} de $A$, et \unsure{vérifient}} \inserted{$N^>(A)$ dép. \unsure{fonctlt}} pour les \uuline{isogénies}, \unsure{ayant|vérifiant} les propriétés suivantes

% slanted
\margin{\deleted{$\xi\in N^+(A)$ : $\exists m>0$ tel que $m\xi$ soit représentable par un diviseur \uline{positif} $\iff$ $\xi$ est l'image inverse d'un \unsure{élt|élément} d'un $N^>(B)$ ; cela \unsure{dispense} donc d'introduire $N^{\unsure{+|>}}(A)$ explicitement.}}

a) $N^>(A)+N^>(A)\subset N^>(A)$, \deleted{\unsure{et}} $m>0\Rightarrow mN^>(A)\subset N^>(A)$
\inserted{Kif Kif pour $N^{\unsure{\geqslant|>}}(A)$}

b) $N^>(A)$ engendre $N(A)$

\deleted{c) $\xi\in N^+(A)\Longrightarrow I$ est $>0$ sur $N^+(A)$, i.e. $\xi_1,\cdots\xi_r\in N^+(A)\Rightarrow I(\xi_1,\cdots\xi_r)>0$}

c) $\xi\in N^>(A)\Longrightarrow\varphi_\xi$ \unsure{$\xi$} est un $\mathbb Q$-isom :

% archive sheet number: 8

